\documentclass[11pt,a4paper]{article}

\usepackage[margin=2.5cm]{geometry}
\usepackage{amsmath,amssymb,amsthm}
\usepackage{bm}
\usepackage{booktabs}
\usepackage{enumitem}
\usepackage[colorlinks=true,linkcolor=blue,citecolor=blue]{hyperref}

\newtheorem{theorem}{Theorem}
\newtheorem{proposition}{Proposition}
\newtheorem{lemma}{Lemma}
\theoremstyle{remark}
\newtheorem{remark}{Remark}

\newcommand{\R}{\mathbb{R}}
\newcommand{\dx}{\,\mathrm{d}x}
\newcommand{\dy}{\,\mathrm{d}y}
\newcommand{\dt}{\,\mathrm{d}t}
\newcommand{\ut}{\tilde{u}}
\newcommand{\us}{u^{*}}
\newcommand{\Om}{\Omega}

\title{\textbf{CE-PINN: Complete Mathematical Derivations}\\[0.4em]
\large Conservation-Embedded Physics-Informed Neural Networks\\
\large for Derivative-Dependent Invariants}
\author{}
\date{}

\begin{document}
\maketitle

\begin{abstract}
\noindent
This document derives every step of the CE-PINN method, from the governing
equations through to the gradient used in backpropagation. Sections marked
\textbf{[PRIOR]} reproduce existing results (PINN-Proj, SDIFP) for completeness
and to establish exactly where they stop. Sections marked \textbf{[NEW]} are the
contribution of this work. Every algebraic step is written out.
\end{abstract}

\tableofcontents

\bigskip\hrule\bigskip

%==========================================================================
\section{Notation}
%==========================================================================

Throughout, $\Om = [a,b] \subset \R$ is the spatial domain with length
$|\Om| = b-a$, and $[0,T]$ is the time interval. The exact solution is $u(x,t)$;
the raw neural network output is $N(x,t;\theta)$ with parameters $\theta$; the
output after imposing boundary conditions is $\ut(x,t)$; and the output after
conservation projection is $\us(x,t)$.

We write $\ut_x = \partial \ut/\partial x$ and so on. Spatial integrals are over
$\Om$ unless stated otherwise. The projection acts on a single time level, so
$t$ is a fixed parameter inside Sections \ref{sec:prior}--\ref{sec:helm} and we
suppress it, writing $\ut(x)$ for $\ut(x,t)$.

For the discrete formulation we use a uniform grid
\begin{equation}
x_i = a + (i-1)\Delta x, \qquad i = 1,\dots,N, \qquad
\Delta x = \frac{b-a}{N-1},
\end{equation}
and identify a function with its grid vector, $\bm{u} = (u_1,\dots,u_N)^\top \in \R^N$.

%==========================================================================
\section{Governing equations and their conservation laws}
%==========================================================================

\subsection{Korteweg--de Vries (KdV)}

\begin{equation}
u_t + 6uu_x + u_{xxx} = 0 .
\label{eq:kdv}
\end{equation}

Conserved integrals:
\begin{align}
C_1 &= \int u \dx && \text{(mass, \emph{linear})}, \label{eq:kdv-c1}\\
C_2 &= \int u^2 \dx && \text{(momentum, \emph{polynomial})}, \label{eq:kdv-c2}\\
C_3 &= \int \Big(u^3 - \tfrac12 u_x^2\Big) \dx && \text{(energy, \emph{derivative-dependent})}. \label{eq:kdv-c3}
\end{align}

One-soliton solution:
\begin{equation}
u(x,t) = \frac{k^2}{2}\,\mathrm{sech}^2\!\left[\frac{k}{2}\big(x - k^2 t\big)\right],
\end{equation}
with amplitude $k^2/2$ and speed $k^2$. On the infinite line,
\begin{equation}
C_1 = 2k, \qquad C_2 = \tfrac{2}{3}k^3, \qquad C_3 = \tfrac{1}{5}k^5 .
\end{equation}

\subsection{Regularised long-wave (RLW)}

\begin{equation}
u_t + u_x + \varepsilon u u_x - \mu u_{xxt} = 0, \qquad \varepsilon>0,\ \mu>0 .
\label{eq:rlw}
\end{equation}

Conserved integrals:
\begin{align}
I_1 &= \int u \dx && \text{(mass, \emph{linear})}, \label{eq:rlw-i1}\\
I_2 &= \int \big(u^2 + \mu u_x^2\big) \dx && \text{(momentum, \emph{derivative-dependent})}, \label{eq:rlw-i2}\\
I_3 &= \int \big(\varepsilon u^3 + 3u^2\big) \dx && \text{(energy, \emph{polynomial})}. \label{eq:rlw-i3}
\end{align}

\begin{remark}
Equation \eqref{eq:rlw-i2} is the central object of this work. The term
$\mu\int u_x^2\dx$ makes $I_2$ \emph{derivative-dependent}. Every existing
projection method is restricted to invariants of the form $\int \phi(u)\dx$ and
therefore cannot enforce $I_2$.
\end{remark}

\begin{remark}[The $\varepsilon$ factor in $I_3$]
Writing \eqref{eq:rlw} as $(1-\mu\partial_x^2)u_t = -\partial_x f(u)$ with
$f(u) = u + \tfrac{\varepsilon}{2}u^2$, and setting
$v = (1-\mu\partial_x^2)^{-1}f$ so that $u_t = -v_x$, gives
\begin{equation}
\frac{\mathrm d}{\mathrm dt}\int\Big(\tfrac12 u^2 + \tfrac{\varepsilon}{6}u^3\Big)\dx
= \int f\,u_t \dx = -\int (v-\mu v_{xx})v_x\dx
= -\tfrac12\big[v^2\big]_a^b + \tfrac{\mu}{2}\big[v_x^2\big]_a^b = 0
\end{equation}
for decaying data. Multiplying by six gives \eqref{eq:rlw-i3}. The factor
$\varepsilon$ is not optional. With $\varepsilon = 1.5$, $\mu = 1/6$, a
pseudospectral integration holds $\int(\varepsilon u^3+3u^2)\dx$ constant to
$10^{-15}$ while $\int(u^3+3u^2)\dx$ drifts by $5\times10^{-4}$ over $t=5$.
The two coincide only at $\varepsilon = 1$.
\end{remark}

Single solitary wave:
\begin{equation}
u(x,t) = 3c\;\mathrm{sech}^2\!\left[\frac{1}{2}\sqrt{\frac{\varepsilon c}{\mu(1+\varepsilon c)}}
\Big(x - (1+\varepsilon c)t - x_0\Big)\right],
\label{eq:rlw-soliton}
\end{equation}
with amplitude $3c$ and speed $1+\varepsilon c$ ($c$ is an amplitude parameter,
not the speed).

\subsection{Kadomtsev--Petviashvili (KP)}

\begin{equation}
\big(u_t + 6uu_x + u_{xxx}\big)_x + \sigma u_{yy} = 0, \qquad \sigma = -3 .
\label{eq:kp}
\end{equation}

Conserved integrals over $\Om_2 = [a_1,b_1]\times[a_2,b_2]$:
\begin{align}
C_1 &= \iint u \dx\dy, \\
C_2 &= \iint u^2 \dx\dy, \\
H &= \iint \Big[-u^3 + \tfrac12 u_x^2 + \tfrac32 \big(\partial_x^{-1}u_y\big)^2\Big]\dx\dy .
\end{align}
The sign of the nonlocal term in $H$ is tied to the convention $\sigma = -3$
used here; KP sign conventions differ across the literature, so $H$ should be
re-derived if \eqref{eq:kp} is written differently.

The lump solution for $\sigma=-3$ is
\begin{equation}
u(x,y,t) = \frac{4\big(-\xi^2 + y^2 + 1\big)}{\big(\xi^2+y^2+1\big)^2},
\qquad \xi = x - 3t,
\end{equation}
for which $C_2 = 8\pi$ exactly.

%==========================================================================
\section{When are these quantities actually conserved?}
%==========================================================================

The integrals above are conserved on the whole line, or on a periodic domain.
On a bounded interval they acquire boundary flux terms, and it is worth being
explicit about this before building a method that pins them to constants.

Integrating \eqref{eq:rlw} over $[a,b]$,
\begin{equation}
\frac{\mathrm d}{\mathrm dt}\int_a^b u \dx
= -\Big[\,u + \tfrac{\varepsilon}{2}u^2 - \mu\,u_{xt}\,\Big]_a^b ,
\label{eq:rlw-flux}
\end{equation}
and for KdV \eqref{eq:kdv},
\begin{equation}
\frac{\mathrm d}{\mathrm dt}\int_a^b u \dx
= -\Big[\,3u^2 + u_{xx}\,\Big]_a^b .
\end{equation}
Analogous boundary terms appear for $I_2$ and $I_3$.

In the implementation these are evaluated from the boundary data alone,
dropping $\mu u_{xt}$ for RLW and $u_{xx}$ for KdV, both of which are
negligible for a solution that has flattened at the domain edge. The check is
therefore a screen for the leading term rather than an exact flux budget; for
non-decaying boundary data it would need the discrete curvature at the ends.

\begin{remark}[Standing assumption]
\label{rem:flux}
Everything that follows assumes the boundary fluxes vanish, which holds for
solitary waves on a domain wide enough that the solution has decayed at both
ends, and for periodic data. Under that assumption $C_1$ and $C_2$ are genuine
constants and projecting onto their initial values is meaningful.
\end{remark}

\begin{remark}[Where the assumption fails]
\label{rem:bore-flux}
It does not hold for the undular bore, whose left boundary is held at
$u_0 = 0.1$ while the right boundary is zero. With a steady left state
($u_{xt}\to0$), \eqref{eq:rlw-flux} gives
\begin{equation}
\frac{\mathrm d}{\mathrm dt}\int u \dx
= u_0 + \tfrac{\varepsilon}{2}u_0^2 \approx 0.108
\qquad (\varepsilon = 1.5),
\end{equation}
against $C_1 \approx 3.6$ on $[-36,300]$. Mass therefore grows by roughly
$150\%$ over $T = 50$: it is injected at the left boundary, and $C_1$ is not
constant. Enforcing the initial $C_1$ there would project onto the wrong
target.

A configuration of this kind needs time-dependent targets $C_1(t)$, $C_2(t)$
obtained by integrating the flux, which is a straightforward extension of the
Newton system (the targets simply move) but is not carried out here. The
implementation reports the flux at set-up rather than silently enforcing a
value that is not conserved.
\end{remark}

%==========================================================================
\section{The PINN baseline \textbf{[PRIOR]}}
%==========================================================================

A physics-informed neural network approximates $u$ by a feed-forward network
$N(x,t;\theta)$ and minimises a composite loss
\begin{equation}
\mathcal{L}(\theta)
= \underbrace{\frac{1}{N_r}\sum_{i=1}^{N_r} \big|r(x_i,t_i)\big|^2}_{\mathcal{L}_{\text{PDE}}}
+ \underbrace{\frac{1}{N_0}\sum_{i=1}^{N_0}\big|N(x_i,0)-f(x_i)\big|^2}_{\mathcal{L}_{\text{IC}}}
+ \underbrace{\frac{1}{N_b}\sum_{i=1}^{N_b}\big|N(x_i^b,t_i)-g(t_i)\big|^2}_{\mathcal{L}_{\text{BC}}},
\end{equation}
where $r$ is the PDE residual, e.g.\ for RLW
\begin{equation}
r = N_t + N_x + \varepsilon N N_x - \mu N_{xxt},
\end{equation}
computed by automatic differentiation.

\begin{remark}[Why soft constraints are insufficient]
$\mathcal{L}_{\text{IC}}$, $\mathcal{L}_{\text{BC}}$ and any conservation
penalty are \emph{soft}: they are minimised, not enforced. Nothing guarantees
they hold after training. Published evidence is strong: replacing the PDE
residual with conservation penalties produced $47.8\%$ error on two-soliton NLS
(Wu et al., 2022). This motivates enforcing physics \emph{by construction}.
\end{remark}

%==========================================================================
\section{Layer 1: hard boundary and initial conditions \textbf{[PRIOR]}}
%==========================================================================

\subsection{Construction}

Consider an evolution equation on $[a,b]\times[0,T]$ with
\begin{equation}
u(x,0) = f(x), \qquad u(a,t) = g_a(t), \qquad u(b,t) = g_b(t),
\end{equation}
subject to the compatibility conditions
\begin{equation}
f(a) = g_a(0), \qquad f(b) = g_b(0).
\label{eq:compat}
\end{equation}

Define, with $s = \dfrac{x-a}{b-a} \in [0,1]$,
\begin{align}
A(x,t) &= f(x) + (1-s)\big[g_a(t)-g_a(0)\big] + s\big[g_b(t)-g_b(0)\big],
\label{eq:A}\\
B(x,t) &= \frac{t}{T}\,\frac{4(x-a)(b-x)}{(b-a)^2}
= \frac{t}{T}\,4s(1-s),
\label{eq:B}
\end{align}
and set
\begin{equation}
\boxed{\ \ut(x,t;\theta) = A(x,t) + B(x,t)\,N(x,t;\theta). \ }
\label{eq:hardbc}
\end{equation}

\begin{remark}[Normalisation of the mask]
The unnormalised choice $B = t(x-a)(b-x)$ has the same zero set and is the
form usually quoted, but its magnitude is $T\big((b-a)/2\big)^2$, which is
about $5\times10^4$ on the RLW domain $[-40,60]$ with $T=20$ and about
$10^6$ for the KP box. The untrained network is then amplified by four to six
orders of magnitude, and the first residual evaluation overflows into the
$10^{14}$ range. Dividing by $T$ and by $\big((b-a)/2\big)^2$ leaves $B$
peaking at $1$ at the domain centre at $t=T$ and changes nothing structural:
$B$ still vanishes at $t=0$ and at both ends.
\end{remark}

\subsection{Verification (all three conditions)}

\paragraph{Initial condition, $t=0$.}
From \eqref{eq:B}, $B(x,0) = 0\cdot(x-a)(b-x) = 0$. From \eqref{eq:A},
\begin{equation}
A(x,0) = f(x) + (1-s)\big[g_a(0)-g_a(0)\big] + s\big[g_b(0)-g_b(0)\big]
= f(x) + 0 + 0 = f(x).
\end{equation}
Hence $\ut(x,0) = f(x) + 0\cdot N = f(x)$. \hfill$\checkmark$

\paragraph{Left boundary, $x=a$.}
Then $s = 0$ and $B(a,t) = t(a-a)(b-a) = 0$. Also
\begin{equation}
A(a,t) = f(a) + 1\cdot\big[g_a(t)-g_a(0)\big] + 0 = g_a(0) + g_a(t) - g_a(0) = g_a(t),
\end{equation}
using $f(a)=g_a(0)$ from \eqref{eq:compat}. Hence $\ut(a,t) = g_a(t)$. \hfill$\checkmark$

\paragraph{Right boundary, $x=b$.}
Then $s=1$ and $B(b,t) = t(b-a)(b-b) = 0$. Also
\begin{equation}
A(b,t) = f(b) + 0 + 1\cdot\big[g_b(t)-g_b(0)\big] = g_b(0)+g_b(t)-g_b(0) = g_b(t).
\end{equation}
Hence $\ut(b,t) = g_b(t)$. \hfill$\checkmark$

All three conditions hold \emph{for every $\theta$}, so $\mathcal{L}_{\text{IC}}$
and $\mathcal{L}_{\text{BC}}$ can be removed from the loss entirely.

\subsection{Two-dimensional version}

On $\Om_2 = [a_1,b_1]\times[a_2,b_2]$ with homogeneous Dirichlet data on all four
edges,
\begin{equation}
\ut(x,y,t;\theta) = f(x,y) + t\,(x-a_1)(b_1-x)(y-a_2)(b_2-y)\,N(x,y,t;\theta).
\end{equation}
At $t=0$ the mask vanishes and $\ut = f$; on each edge one factor vanishes.

%==========================================================================
\section{Layer 2, prior work: projections for polynomial invariants \textbf{[PRIOR]}}
\label{sec:prior}
%==========================================================================

All projections in this section solve
\begin{equation}
\min_{\us} \ \int_a^b \big(\us - \ut\big)^2 \dx
\quad\text{subject to prescribed integral constraints,}
\end{equation}
i.e.\ they find the closest function (in $L^2$) satisfying the conservation laws.

\subsection{Case A: single linear constraint (PINN-Proj, Eq.~12)}

Constraint: $\displaystyle\int \us \dx = C_1$.

\paragraph{Step 1. Lagrangian.}
\begin{equation}
\mathcal{J}[\us,\lambda] = \int (\us-\ut)^2\dx + \lambda\left(\int \us \dx - C_1\right).
\end{equation}

\paragraph{Step 2. Variation.}
Perturb $\us \to \us + \epsilon v$ and differentiate at $\epsilon = 0$:
\begin{equation}
\frac{\mathrm{d}}{\mathrm{d}\epsilon}\bigg|_{\epsilon=0}
\int (\us+\epsilon v-\ut)^2\dx = \int 2(\us-\ut)\,v \dx,
\qquad
\frac{\mathrm{d}}{\mathrm{d}\epsilon}\bigg|_{\epsilon=0}
\lambda\int(\us+\epsilon v)\dx = \lambda\int v\dx .
\end{equation}
Setting the total variation to zero for all $v$:
\begin{equation}
\int\Big[2(\us-\ut) + \lambda\Big]v\dx = 0 \quad \forall v
\qquad\Longrightarrow\qquad
2(\us-\ut) + \lambda = 0 .
\end{equation}

\paragraph{Step 3. Solve pointwise.}
\begin{equation}
\us = \ut - \frac{\lambda}{2}.
\label{eq:lin-form}
\end{equation}

\paragraph{Step 4. Impose the constraint.}
Substituting \eqref{eq:lin-form}:
\begin{equation}
\int\left(\ut - \frac{\lambda}{2}\right)\dx = C_1
\quad\Longrightarrow\quad
\int \ut \dx - \frac{\lambda}{2}\,|\Om| = C_1
\quad\Longrightarrow\quad
\frac{\lambda}{2} = \frac{1}{|\Om|}\left(\int\ut\dx - C_1\right).
\end{equation}

\paragraph{Step 5. Result.}
\begin{equation}
\boxed{\ \us(x) = \ut(x) - \frac{1}{|\Om|}\int \ut \dx + \frac{C_1}{|\Om|}. \ }
\label{eq:proj-linear}
\end{equation}
This is the de-mean-and-reshift correction of PINN-Proj. It is closed-form and
costs $O(N)$.

\subsection{Case B: single quadratic constraint (PINN-Proj, Eq.~15)}

Constraint: $\displaystyle\int (\us)^2 \dx = C_2$.

\paragraph{Step 1--2. Lagrangian and variation.}
\begin{equation}
\mathcal{J} = \int(\us-\ut)^2\dx + \lambda\left(\int(\us)^2\dx - C_2\right),
\end{equation}
\begin{equation}
\delta\mathcal{J} = \int\Big[2(\us-\ut) + 2\lambda \us\Big]v\dx = 0
\quad\Longrightarrow\quad
2(\us-\ut) + 2\lambda\us = 0 .
\end{equation}

\paragraph{Step 3. Solve.}
\begin{equation}
\us(1+\lambda) = \ut
\qquad\Longrightarrow\qquad
\us = \frac{\ut}{1+\lambda}.
\end{equation}

\paragraph{Step 4. Impose the constraint.}
\begin{equation}
\int \frac{\ut^2}{(1+\lambda)^2}\dx = C_2
\quad\Longrightarrow\quad
(1+\lambda)^2 = \frac{\int \ut^2\dx}{C_2}
\quad\Longrightarrow\quad
1+\lambda = \sqrt{\frac{\int\ut^2\dx}{C_2}} .
\end{equation}

\paragraph{Step 5. Result.}
\begin{equation}
\boxed{\ \us(x) = \ut(x)\,\sqrt{\frac{C_2}{\int \ut^2 \dx}}. \ }
\label{eq:proj-quad}
\end{equation}
A pure rescaling. Note this destroys $\int\us\dx = C_1$ unless $C_1 = 0$.

\subsection{Case C: both constraints jointly (PINN-Proj Eq.~16--17; SDIFP Eq.~4.2)}

Constraints: $\displaystyle\int\us\dx = C_1$ \emph{and} $\displaystyle\int(\us)^2\dx = C_2$.

\paragraph{Step 1--2. Lagrangian and variation.}
\begin{equation}
\mathcal{J} = \int(\us-\ut)^2\dx
+ \lambda_1\left(\int\us\dx - C_1\right)
+ \lambda_2\left(\int(\us)^2\dx - C_2\right),
\end{equation}
\begin{equation}
\delta\mathcal{J}
= \int\Big[2(\us-\ut) + \lambda_1 + 2\lambda_2\us\Big]v\dx = 0
\quad\Longrightarrow\quad
2(\us-\ut) + \lambda_1 + 2\lambda_2\us = 0 .
\label{eq:joint-el}
\end{equation}

\paragraph{Step 3. Solve pointwise.}
Collecting $\us$ terms in \eqref{eq:joint-el}:
\begin{equation}
2\us + 2\lambda_2\us = 2\ut - \lambda_1
\quad\Longrightarrow\quad
\us(1+\lambda_2) = \ut - \frac{\lambda_1}{2}
\quad\Longrightarrow\quad
\us = \frac{\ut - \lambda_1/2}{1+\lambda_2}.
\label{eq:joint-form}
\end{equation}

\paragraph{Step 4. Recognise the affine structure.}
Define
\begin{equation}
\alpha = \frac{1}{1+\lambda_2}, \qquad
\beta = \frac{-\lambda_1/2}{1+\lambda_2},
\end{equation}
so that \eqref{eq:joint-form} reads
\begin{equation}
\us(x) = \alpha\,\ut(x) + \beta .
\label{eq:affine}
\end{equation}
\emph{The optimal projection is an affine map of the network output.} This is the
key structural fact about all prior work.

\paragraph{Step 5. Determine $\alpha,\beta$ from the constraints.}
Introduce the moments
\begin{equation}
\tilde m = \frac{1}{|\Om|}\int \ut \dx, \qquad
\tilde s = \frac{1}{|\Om|}\int \ut^2 \dx, \qquad
\tilde\sigma^2 = \tilde s - \tilde m^2,
\end{equation}
and the targets
\begin{equation}
m_0 = \frac{C_1}{|\Om|}, \qquad
s_0 = \frac{C_2}{|\Om|}, \qquad
\sigma_0^2 = s_0 - m_0^2 .
\end{equation}

\emph{First constraint.} Integrating \eqref{eq:affine},
\begin{equation}
\int\us\dx = \alpha\int\ut\dx + \beta|\Om| = C_1
\quad\Longrightarrow\quad
\alpha\tilde m + \beta = m_0 .
\label{eq:c1-moment}
\end{equation}

\emph{Second constraint.} Squaring \eqref{eq:affine} and integrating,
\begin{equation}
\int(\us)^2\dx = \alpha^2\int\ut^2\dx + 2\alpha\beta\int\ut\dx + \beta^2|\Om| = C_2,
\end{equation}
and dividing by $|\Om|$,
\begin{equation}
\alpha^2\tilde s + 2\alpha\beta\tilde m + \beta^2 = s_0 .
\label{eq:c2-moment}
\end{equation}

\emph{Eliminate $\beta$.} From \eqref{eq:c1-moment}, $\beta = m_0 - \alpha\tilde m$.
Substituting into \eqref{eq:c2-moment}:
\begin{align}
s_0 &= \alpha^2\tilde s + 2\alpha(m_0-\alpha\tilde m)\tilde m + (m_0-\alpha\tilde m)^2 \\
&= \alpha^2\tilde s + 2\alpha m_0\tilde m - 2\alpha^2\tilde m^2
   + m_0^2 - 2\alpha m_0\tilde m + \alpha^2\tilde m^2 \\
&= \alpha^2\tilde s - \alpha^2\tilde m^2 + m_0^2 \\
&= \alpha^2\big(\tilde s - \tilde m^2\big) + m_0^2
 = \alpha^2\tilde\sigma^2 + m_0^2 .
\end{align}
Therefore
\begin{equation}
\alpha^2\tilde\sigma^2 = s_0 - m_0^2 = \sigma_0^2
\quad\Longrightarrow\quad
\alpha = \frac{\sigma_0}{\tilde\sigma},
\qquad
\beta = m_0 - \frac{\sigma_0}{\tilde\sigma}\tilde m .
\end{equation}

\paragraph{Step 6. Result.}
\begin{equation}
\boxed{\ \us(x) = \frac{\sigma_0}{\tilde\sigma}\Big(\ut(x) - \tilde m\Big) + m_0 . \ }
\label{eq:proj-joint}
\end{equation}
This is exactly the SDIFP standardisation transform. It is closed-form, $O(N)$,
and matches the first two moments of the target.

\begin{remark}[Degenerate cases]
The closed forms above assume non-degenerate moments. Equation
\eqref{eq:proj-quad} requires $\int\ut^2\dx > 0$ and $C_2 > 0$, and fixes the
sign of the root by taking the positive branch, which keeps $\us$ aligned with
$\ut$. Equation \eqref{eq:proj-joint} requires $\tilde\sigma > 0$, so $\ut$
must not be constant, and $\sigma_0^2 = C_2/|\Om| - m_0^2 \geq 0$, which is the
feasibility condition $C_2 \geq C_1^2/|\Om|$ (Cauchy--Schwarz, with equality
only for constant $\us$). If $\tilde\sigma = 0$ but $\sigma_0 > 0$ the target is
simply unreachable: a constant network output cannot be rescaled into one with
positive variance. If $\tilde\sigma = \sigma_0 = 0$ the map is $0/0$ and should
be read as the identity when the mass constraint already holds, and as the
constant $m_0$ otherwise. Implementations clamp both variances away from zero.
\end{remark}

%==========================================================================
\section{Why prior work cannot handle derivative-dependent invariants \textbf{[NEW analysis]}}
%==========================================================================

We now make precise the limitation that both PINN-Proj and SDIFP state in words.

Consider the RLW momentum \eqref{eq:rlw-i2},
\begin{equation}
I_2[\us] = \int \Big[(\us)^2 + \mu(\us_x)^2\Big]\dx = C_2 .
\end{equation}

\begin{proposition}[The optimal projection is not affine]
\label{prop:not-affine}
Let $\mu > 0$, and let $1$, $\ut$ and $\ut_{xx}$ be linearly independent on
$\Om$. Then the minimiser of $\int(\us-\ut)^2\dx$ subject to
$\int\us\dx = C_1$ and $\int[(\us)^2+\mu(\us_x)^2]\dx = C_2$ has the form
$\us = \alpha\ut + \beta$ with constants $\alpha,\beta$ \emph{only if}
$\lambda_2 = 0$.

All three functions are needed in the hypothesis, because the substitution
below leaves a constant term alongside $\ut$ and $\ut_{xx}$; independence of
the pair alone would not isolate the $\ut_{xx}$ coefficient. The hypothesis is
generic but not automatic, and fails exactly when $\ut_{xx} = a\ut + b$ for
constants $a,b$, in which case an affine minimiser can occur with
$\lambda_2 \neq 0$.

We state the conclusion as $\lambda_2 = 0$ rather than ``the constraint is
inactive''. An equality constraint can be satisfied and still carry a zero
multiplier; the operative fact is that the elliptic term disappears.
\end{proposition}

\begin{proof}
By Theorem \ref{thm:el} below, the minimiser satisfies
\begin{equation}
(1+\lambda_2)\us - \lambda_2\mu\,\us_{xx} = \ut - \frac{\lambda_1}{2}.
\end{equation}
Suppose $\us = \alpha\ut+\beta$. Then $\us_{xx} = \alpha\ut_{xx}$ and substitution gives
\begin{equation}
(1+\lambda_2)\big(\alpha\ut+\beta\big) - \lambda_2\mu\,\alpha\ut_{xx}
= \ut - \frac{\lambda_1}{2},
\end{equation}
that is,
\begin{equation}
\underbrace{\Big[\alpha(1+\lambda_2)-1\Big]}_{\text{coefficient of }\ut}\ut
\;-\;\underbrace{\lambda_2\mu\alpha}_{\text{coefficient of }\ut_{xx}}\ut_{xx}
\;+\;\Big[(1+\lambda_2)\beta + \tfrac{\lambda_1}{2}\Big] = 0 .
\end{equation}
For this to hold for a general profile $\ut$, in which $\ut$ and $\ut_{xx}$ are
linearly independent, each coefficient must vanish. The second requires
$\lambda_2\mu\alpha = 0$. Since $\mu>0$ and $\alpha\neq 0$ (else $\us$ is
constant), we need $\lambda_2 = 0$, i.e.\ the momentum constraint carries no
multiplier and is inactive.
\end{proof}

\begin{remark}[Two distinct failures]
Proposition \ref{prop:not-affine} identifies the essential obstruction, and there
are two practical consequences.

\emph{(i) Suboptimality.} The affine family $\{\alpha\ut+\beta\}$ has only two
degrees of freedom. Two constraints therefore determine it uniquely, leaving no
freedom to minimise $\|\us-\ut\|$. One can force an affine map to hit the
\emph{numerical value} of $C_2$, but the resulting $\us$ is not the closest point
on the constraint manifold, so it perturbs the solution more than necessary.

\emph{(ii) Boundary conditions.} From \eqref{eq:affine},
$\us(a) = \alpha g_a + \beta$ and $\us(b) = \alpha g_b + \beta$. Preserving both
requires $\alpha g_a + \beta = g_a$ and $\alpha g_b + \beta = g_b$. For
$g_a \neq g_b$ these force $\alpha = 1$, $\beta = 0$, so only the identity
preserves the data. For $g_a = g_b = g$ they collapse to the single condition
$\beta = (1-\alpha)g$, which \emph{can} hold with $\alpha \neq 1$; in
particular any map with $\beta = 0$ preserves homogeneous data. The affine
projection sets $\beta = m_0 - \alpha\tilde m$, which satisfies neither
condition in general, so non-homogeneous Dirichlet data is destroyed. This is
the content of SDIFP Remark 4.6.
\end{remark}

%==========================================================================
\section{Layer 2, this work: the Helmholtz projection \textbf{[NEW]}}
\label{sec:helm}
%==========================================================================

\subsection{Problem statement}

Given the boundary-condition-satisfying output $\ut$, find
\begin{equation}
\us = \arg\min_{w} \int_a^b \big(w-\ut\big)^2\dx
\label{eq:helm-obj}
\end{equation}
subject to
\begin{align}
G_1[w] &:= \int_a^b w \dx - C_1 = 0, \label{eq:helm-g1}\\
G_2[w] &:= \int_a^b \Big[w^2 + \mu (w_x)^2\Big]\dx - C_2 = 0, \label{eq:helm-g2}\\
w(a) &= g_a, \qquad w(b) = g_b. \label{eq:helm-bc}
\end{align}

\subsection{Derivation of the Euler--Lagrange equation}

\begin{theorem}[Helmholtz Euler--Lagrange equation]
\label{thm:el}
The minimiser of \eqref{eq:helm-obj}--\eqref{eq:helm-bc} satisfies
\begin{equation}
\boxed{\ (1+\lambda_2)\,\us - \lambda_2\mu\,\us_{xx} = \ut - \frac{\lambda_1}{2},
\qquad \us(a)=g_a,\ \us(b)=g_b, \ }
\label{eq:helmholtz}
\end{equation}
for scalars $\lambda_1,\lambda_2$ determined by \eqref{eq:helm-g1}--\eqref{eq:helm-g2}.
\end{theorem}

\begin{proof}
Form the Lagrangian
\begin{equation}
\mathcal{J}[\us,\lambda_1,\lambda_2]
= \int(\us-\ut)^2\dx
+ \lambda_1\left(\int\us\dx - C_1\right)
+ \lambda_2\left(\int\big[(\us)^2 + \mu(\us_x)^2\big]\dx - C_2\right).
\end{equation}

Perturb $\us \to \us + \epsilon v$ where $v$ is any smooth function with
\begin{equation}
v(a) = v(b) = 0
\label{eq:vbc}
\end{equation}
(the boundary values of $\us$ are prescribed, so admissible variations vanish
there). We compute $\frac{\mathrm d}{\mathrm d\epsilon}\mathcal{J}\big|_{\epsilon=0}$
term by term.

\paragraph{Term 1: the objective.}
\begin{equation}
\frac{\mathrm d}{\mathrm d\epsilon}\int(\us+\epsilon v-\ut)^2\dx
= \int 2(\us+\epsilon v-\ut)\,v\dx
\;\xrightarrow[\epsilon\to0]{}\;
\int 2(\us-\ut)\,v\dx .
\end{equation}

\paragraph{Term 2: the mass constraint.}
\begin{equation}
\frac{\mathrm d}{\mathrm d\epsilon}\,\lambda_1\int(\us+\epsilon v)\dx
= \lambda_1\int v\dx .
\end{equation}

\paragraph{Term 3: the momentum constraint.}
Write the integrand as $(\us+\epsilon v)^2 + \mu(\us_x+\epsilon v_x)^2$. Then
\begin{equation}
\frac{\mathrm d}{\mathrm d\epsilon}\,\lambda_2\int\Big[(\us+\epsilon v)^2
+ \mu(\us_x+\epsilon v_x)^2\Big]\dx
= \lambda_2\int\Big[2(\us+\epsilon v)v + 2\mu(\us_x+\epsilon v_x)v_x\Big]\dx,
\end{equation}
which at $\epsilon=0$ becomes
\begin{equation}
\lambda_2\int 2\us v\dx + \lambda_2\int 2\mu \us_x v_x \dx .
\label{eq:term3}
\end{equation}

The second integral in \eqref{eq:term3} contains $v_x$, not $v$, so we integrate
by parts:
\begin{equation}
\int_a^b 2\mu\,\us_x\,v_x\dx
= \Big[2\mu\,\us_x\,v\Big]_a^b - \int_a^b 2\mu\,\us_{xx}\,v\dx .
\end{equation}
By \eqref{eq:vbc} the boundary term vanishes:
$\big[2\mu\us_x v\big]_a^b = 2\mu\us_x(b)v(b) - 2\mu\us_x(a)v(a) = 0 - 0 = 0$.
Hence
\begin{equation}
\int 2\mu\,\us_x v_x\dx = -\int 2\mu\,\us_{xx}\,v\dx .
\label{eq:ibp}
\end{equation}

\paragraph{Assemble.} Collecting all terms,
\begin{equation}
\delta\mathcal{J}
= \int\Big[\,2(\us-\ut) + \lambda_1 + 2\lambda_2\us - 2\lambda_2\mu\,\us_{xx}\,\Big]\,v\dx = 0
\qquad\text{for all admissible } v .
\end{equation}
By the fundamental lemma of the calculus of variations, the bracket vanishes
pointwise:
\begin{equation}
2(\us-\ut) + \lambda_1 + 2\lambda_2\us - 2\lambda_2\mu\,\us_{xx} = 0 .
\label{eq:el-raw}
\end{equation}

\paragraph{Rearrange.} Divide \eqref{eq:el-raw} by $2$:
\begin{equation}
\us - \ut + \frac{\lambda_1}{2} + \lambda_2\us - \lambda_2\mu\,\us_{xx} = 0 .
\end{equation}
Move $\ut$ and $\lambda_1/2$ to the right and collect $\us$:
\begin{equation}
(1+\lambda_2)\us - \lambda_2\mu\,\us_{xx} = \ut - \frac{\lambda_1}{2}. \qedhere
\end{equation}
\end{proof}

\subsection{Structure of the resulting operator}

Define the \emph{Helmholtz operator}
\begin{equation}
\mathcal{H} := (1+\lambda_2)\,\mathcal{I} - \lambda_2\mu\,\frac{\mathrm d^2}{\mathrm dx^2},
\end{equation}
so that \eqref{eq:helmholtz} reads $\mathcal{H}\us = \ut - \lambda_1/2$.

Dividing through by $(1+\lambda_2)$, assuming $\lambda_2 \neq -1$, gives the
normalised form
\begin{equation}
\us - \kappa\,\us_{xx} = h,
\qquad
\kappa := \frac{\lambda_2\mu}{1+\lambda_2},
\qquad
h := \frac{\ut - \lambda_1/2}{1+\lambda_2}.
\label{eq:helm-normalised}
\end{equation}

\begin{proposition}[Prior work is the limiting case $\lambda_2=0$]
\label{prop:limit}
Setting $\lambda_2 = 0$ in \eqref{eq:helmholtz} gives $\kappa=0$ and
\begin{equation}
\us = \ut - \frac{\lambda_1}{2},
\end{equation}
which is precisely the PINN-Proj linear projection \eqref{eq:lin-form}. More
generally, if the momentum constraint contains no derivative term
($\mu \to 0$), then $\kappa \to 0$ and \eqref{eq:helm-normalised} reduces to
$\us = (\ut-\lambda_1/2)/(1+\lambda_2)$, the affine map \eqref{eq:joint-form} of
prior work.
\end{proposition}

\noindent
Proposition \ref{prop:limit} states the relationship of this work to the
literature: existing projections are the degenerate case in which the elliptic
term is absent. The term $\lambda_2\mu\,\us_{xx}$ is exactly the correction
required by the derivative-dependent part of the invariant, and it acts as a
smoothing operator that redistributes the solution so that the $\int u_x^2$
content matches the target.

\begin{proposition}[Well-posedness]
For $\kappa > 0$ the operator $\mathcal{A} := \mathcal{I} - \kappa\,\partial_x^2$
with Dirichlet boundary conditions is self-adjoint and positive definite, so
\eqref{eq:helm-normalised} has a unique solution.
\end{proposition}

\begin{proof}
On $[a,b]$ with homogeneous Dirichlet data, $-\partial_x^2$ has eigenfunctions
$\sin\!\big(n\pi(x-a)/|\Om|\big)$ with eigenvalues $\big(n\pi/|\Om|\big)^2 > 0$.
Hence $\mathcal{A}$ has eigenvalues $1 + \kappa(n\pi/|\Om|)^2 \geq 1 > 0$.
\end{proof}

\begin{proposition}[The projection is well posed only for $\lambda_2 \geq 0$]
\label{prop:wellposed}
The second variation of $\mathcal{J}$ at a stationary point, taken over
admissible $v$ with $v(a)=v(b)=0$, is
\begin{equation}
\delta^2\mathcal{J}[v]
= 2\int v^2\dx + 2\lambda_2\int\big(v^2 + \mu v_x^2\big)\dx
= 2\int\Big[(1+\lambda_2)v^2 + \lambda_2\mu\,v_x^2\Big]\dx .
\end{equation}
For $\lambda_2 \geq 0$ this is non-negative and the stationary point is a
minimum. For $\lambda_2 < 0$ it can be made negative by taking $v$ with
$\int v_x^2$ large relative to $\int v^2$, i.e.\ a high-frequency perturbation,
so the stationary point is not a minimum.
\end{proposition}

\begin{remark}[What goes wrong, and what to do about it]
\label{rem:lam2neg}
The sign of $\lambda_2$ is set by the data: $\lambda_2 > 0$ when
$I_2[\ut] > C_2$ and the projection must \emph{reduce} the derivative content,
$\lambda_2 < 0$ when it must \emph{increase} it.

In the second case the continuous problem may have no minimiser at all. One can
add oscillations of vanishing amplitude and growing frequency which raise
$\int u_x^2$ to any prescribed level while changing $\ut$ arbitrarily little in
$L^2$, so the infimum is approached but not attained. Consistently with
Proposition \ref{prop:wellposed}, the operator
$\mathcal{I}-\kappa\partial_x^2$ with $\kappa<0$ is
$\mathcal{I}+|\kappa|\partial_x^2$, which is indefinite.

On a grid a minimiser exists, because the mesh caps the available frequency,
but the problem is ill conditioned. The discrete Laplacian has eigenvalues down
to $-4/\Delta x^2$, so positive definiteness survives only while
$|\kappa| < \Delta x^2/4$, that is
\begin{equation}
\lambda_2 \;\geq\; \lambda_2^{\min}
:= \frac{\kappa_{\min}}{\mu - \kappa_{\min}},
\qquad \kappa_{\min} = -\tfrac{1}{4}\Delta x^2 .
\label{eq:lam2min}
\end{equation}
Two conditions must be kept apart here. The continuous second-order optimality
condition is $\lambda_2 \geq 0$: only on that branch is the stationary point a
genuine minimiser. The bound \eqref{eq:lam2min} is weaker and permits a small
negative range; it is a \emph{discrete solvability} condition, the largest
$|\kappa|$ for which the grid operator stays positive definite, and a solution
obtained there should be read as a numerical fallback rather than as a
continuous minimiser.

The implementation clamps $\lambda_2$ at \eqref{eq:lam2min} (with a safety
factor) and records when the clamp binds, which is the honest reading: the
target is not attainable by an elliptic projection on that grid. In practice
this is rare, because an undertrained network carries excess roughness and
therefore excess $I_2$, putting it on the well-posed branch $\lambda_2 > 0$.
The Helmholtz projection is best understood as a \emph{smoothing} projection.

Note also that $\lambda_2 = -1$ must be excluded outright, since the normalised
form \eqref{eq:helm-normalised} divides by $1+\lambda_2$.
\end{remark}

\begin{remark}[Boundary conditions are preserved exactly]
Unlike the affine projection, \eqref{eq:helmholtz} carries Dirichlet data as part
of the boundary value problem. The projected solution satisfies
$\us(a)=g_a$ and $\us(b)=g_b$ by construction, so Layer 1 and Layer 2 compose
without conflict. This is a structural advantage over standardisation, which
shifts boundary values (SDIFP Remark 4.6).
\end{remark}

%==========================================================================
\section{Discretisation \textbf{[NEW]}}
%==========================================================================

\subsection{The tridiagonal system}

Approximate the second derivative by the standard three-point stencil
\begin{equation}
\us_{xx}(x_i) \approx \frac{\us_{i-1} - 2\us_i + \us_{i+1}}{\Delta x^2},
\qquad i = 2,\dots,N-1 .
\end{equation}
Substituting into $\us - \kappa\us_{xx} = h$ at an interior node,
\begin{equation}
\us_i - \frac{\kappa}{\Delta x^2}\Big(\us_{i-1} - 2\us_i + \us_{i+1}\Big) = h_i,
\end{equation}
and collecting terms with $c := \kappa/\Delta x^2$,
\begin{equation}
-c\,\us_{i-1} + \big(1+2c\big)\us_i - c\,\us_{i+1} = h_i,
\qquad i = 2,\dots,N-1 .
\label{eq:interior-row}
\end{equation}
The boundary nodes are prescribed:
\begin{equation}
\us_1 = g_a, \qquad \us_N = g_b .
\label{eq:boundary-row}
\end{equation}

In matrix form $\bm{A}\bm{\us} = \bm{h}$ with
\begin{equation}
\bm{A} =
\begin{pmatrix}
1 & 0 & & & \\
-c & 1+2c & -c & & \\
 & \ddots & \ddots & \ddots & \\
 & & -c & 1+2c & -c \\
 & & & 0 & 1
\end{pmatrix},
\qquad
\bm{h} =
\begin{pmatrix}
g_a \\ h_2 \\ \vdots \\ h_{N-1} \\ g_b
\end{pmatrix}.
\label{eq:matrixA}
\end{equation}

\begin{remark}[Implementation warning]
\label{rem:bcwarning}
Only the \emph{super}-diagonal entry of row $1$ and the \emph{sub}-diagonal entry
of row $N$ are set to zero. The sub-diagonal entry of row $2$ (the coefficient
$-c$ multiplying $\us_1$) and the super-diagonal entry of row $N-1$ must be
retained: they carry the boundary data into the interior. Zeroing them decouples
the first and last interior nodes from the boundary. When $g_a=g_b=0$ this
happens to be harmless because the dropped terms multiply zero, but it is
incorrect for non-homogeneous data such as the undular bore, where $g_a = u_0 \neq 0$.
\end{remark}

\subsection{The Thomas algorithm}

For a tridiagonal system with sub-diagonal $\{a_i\}_{i=2}^{N}$, diagonal
$\{b_i\}_{i=1}^{N}$, super-diagonal $\{c_i\}_{i=1}^{N-1}$ and right-hand side
$\{d_i\}$:

\paragraph{Forward elimination.}
\begin{align}
c'_1 &= \frac{c_1}{b_1}, &  d'_1 &= \frac{d_1}{b_1}, \\
c'_i &= \frac{c_i}{b_i - a_i c'_{i-1}}\ \ (i \le N-1), &
d'_i &= \frac{d_i - a_i d'_{i-1}}{b_i - a_i c'_{i-1}},
\qquad i = 2,\dots,N .
\end{align}

\paragraph{Back substitution.}
\begin{equation}
\us_N = d'_N, \qquad
\us_i = d'_i - c'_i\,\us_{i+1}, \qquad i = N-1,\dots,1 .
\end{equation}

Total cost: $O(N)$ arithmetic operations, with no fill-in.

\begin{remark}[Vectorised alternative]
The Thomas recursion is inherently sequential. On a GPU, executing $2N$
dependent steps as separate kernel launches dominates the runtime. For the grid
sizes used here ($N \leq 1001$) it is faster in practice to assemble
\eqref{eq:matrixA} densely and call a batched LU solve for all right-hand sides
simultaneously. The arithmetic cost rises from $O(N)$ to $O(N^3)$ but the wall
clock falls by one to two orders of magnitude, because a single dense solve is
one kernel rather than thousands.
\end{remark}

\subsection{Quadrature and discrete derivative}

Let $\bm{q}\in\R^N$ be the trapezoidal weight vector,
\begin{equation}
q_1 = q_N = \frac{\Delta x}{2}, \qquad q_i = \Delta x \ \ (2\le i\le N-1),
\end{equation}
so that $\int w\dx \approx \bm{q}^\top\bm{w}$. Let $\bm{D}\in\R^{N\times N}$ be
the central-difference first-derivative matrix,
\begin{equation}
(\bm{D}\bm{w})_i = \frac{w_{i+1}-w_{i-1}}{2\Delta x}\ \ (2\le i\le N-1),
\qquad
(\bm{D}\bm{w})_1 = \frac{-3w_1+4w_2-w_3}{2\Delta x},
\qquad
(\bm{D}\bm{w})_N = \frac{3w_N-4w_{N-1}+w_{N-2}}{2\Delta x},
\end{equation}
so $\bm{D}$ is second-order central in the interior and second-order one-sided
at the two end nodes.
The discrete constraints are then
\begin{equation}
F_1(\lambda_1,\lambda_2) = \bm{q}^\top\bm{\us} - C_1,
\qquad
F_2(\lambda_1,\lambda_2) = \bm{q}^\top\Big[\bm{\us}\odot\bm{\us}
+ \mu\,(\bm{D}\bm{\us})\odot(\bm{D}\bm{\us})\Big] - C_2,
\label{eq:discrete-F}
\end{equation}
where $\odot$ denotes elementwise product.

\begin{remark}[Consistency of the discrete objects]
Three discrete choices have to agree: the quadrature $\bm q$, the derivative
$\bm{D}$ inside the invariant, and the stencil inside the Helmholtz matrix. The
reference value $C_2$ must be computed with the same $\bm q$ and $\bm{D}$ used
in the constraint, otherwise the discrete target is not attainable and Newton
converges to a residual floor rather than to zero.

The pairing here is second-order throughout: the three-point Laplacian in
$\bm{A}$, the matching central difference in $F_2$, and trapezoidal weights.
This is a consistent finite-difference discretisation of the continuous
projection, and Newton drives $F_1 = F_2 = 0$ exactly, so the \emph{discrete}
invariants are conserved to machine precision. It is not, however, an exactly
structure-preserving discrete variational projection: that would call for
summation-by-parts operators satisfying a discrete integration-by-parts
identity compatible with $\bm q$. The higher-order operators used for the PDE
residual are deliberately kept separate from these.
\end{remark}

%==========================================================================
\section{The Newton system for the multipliers \textbf{[NEW]}}
%==========================================================================

Equation \eqref{eq:helmholtz} determines $\us$ once $(\lambda_1,\lambda_2)$ are
known. The multipliers are fixed by requiring $F_1 = F_2 = 0$. This is a
$2\times 2$ nonlinear system, solved by Newton's method:
\begin{equation}
\begin{pmatrix}\lambda_1\\\lambda_2\end{pmatrix}^{(k+1)}
=
\begin{pmatrix}\lambda_1\\\lambda_2\end{pmatrix}^{(k)}
- \bm{J}^{-1}
\begin{pmatrix}F_1\\F_2\end{pmatrix}^{(k)},
\qquad
\bm{J} = \begin{pmatrix}
\partial F_1/\partial\lambda_1 & \partial F_1/\partial\lambda_2 \\
\partial F_2/\partial\lambda_1 & \partial F_2/\partial\lambda_2
\end{pmatrix}.
\end{equation}

\subsection{Sensitivity equations}

Define $p_1 := \partial\us/\partial\lambda_1$ and
$p_2 := \partial\us/\partial\lambda_2$.

\paragraph{Sensitivity with respect to $\lambda_1$.}
Differentiate \eqref{eq:helmholtz} with respect to $\lambda_1$, holding
$\lambda_2$ and $\ut$ fixed. The left-hand side gives
\begin{equation}
\frac{\partial}{\partial\lambda_1}\Big[(1+\lambda_2)\us - \lambda_2\mu\us_{xx}\Big]
= (1+\lambda_2)p_1 - \lambda_2\mu\,(p_1)_{xx},
\end{equation}
and the right-hand side gives
\begin{equation}
\frac{\partial}{\partial\lambda_1}\left[\ut - \frac{\lambda_1}{2}\right] = -\frac12 .
\end{equation}
Therefore
\begin{equation}
\boxed{\ \mathcal{H}p_1 = -\tfrac12, \qquad p_1(a) = p_1(b) = 0. \ }
\label{eq:p1}
\end{equation}
The homogeneous boundary conditions follow because $\us(a) = g_a$ is prescribed
independently of $\lambda_1$, so $\partial\us(a)/\partial\lambda_1 = 0$.

\paragraph{Sensitivity with respect to $\lambda_2$.}
Differentiate \eqref{eq:helmholtz} with respect to $\lambda_2$. Using the product
rule on both terms of the left-hand side,
\begin{equation}
\frac{\partial}{\partial\lambda_2}\Big[(1+\lambda_2)\us\Big]
= \us + (1+\lambda_2)p_2,
\qquad
\frac{\partial}{\partial\lambda_2}\Big[\lambda_2\mu\,\us_{xx}\Big]
= \mu\,\us_{xx} + \lambda_2\mu\,(p_2)_{xx}.
\end{equation}
The right-hand side is independent of $\lambda_2$, so its derivative is zero.
Hence
\begin{equation}
\us + (1+\lambda_2)p_2 - \mu\,\us_{xx} - \lambda_2\mu\,(p_2)_{xx} = 0,
\end{equation}
that is,
\begin{equation}
\boxed{\ \mathcal{H}p_2 = \mu\,\us_{xx} - \us, \qquad p_2(a) = p_2(b) = 0. \ }
\label{eq:p2}
\end{equation}

\begin{remark}
Both sensitivity problems use \emph{the same operator} $\mathcal{H}$ as the
forward solve. One factorisation therefore serves three right-hand sides:
$\ut - \lambda_1/2$, $-\tfrac12$, and $\mu\us_{xx}-\us$.
\end{remark}

\subsection{Jacobian entries}

\paragraph{First row.} Since $F_1 = \int\us\dx - C_1$,
\begin{equation}
J_{11} = \frac{\partial F_1}{\partial\lambda_1}
= \int \frac{\partial \us}{\partial\lambda_1}\dx = \int p_1 \dx,
\qquad
J_{12} = \frac{\partial F_1}{\partial\lambda_2} = \int p_2 \dx .
\end{equation}

\paragraph{Second row.} Since $F_2 = \int\big[(\us)^2 + \mu(\us_x)^2\big]\dx - C_2$,
apply the chain rule inside the integral:
\begin{equation}
\frac{\partial}{\partial\lambda_j}\Big[(\us)^2 + \mu(\us_x)^2\Big]
= 2\us\,\frac{\partial\us}{\partial\lambda_j}
+ 2\mu\,\us_x\,\frac{\partial\us_x}{\partial\lambda_j}
= 2\us p_j + 2\mu\,\us_x\,(p_j)_x ,
\end{equation}
where we used $\partial(\us_x)/\partial\lambda_j = (\partial\us/\partial\lambda_j)_x = (p_j)_x$.
Hence
\begin{equation}
J_{21} = \int\Big[2\us p_1 + 2\mu\,\us_x (p_1)_x\Big]\dx,
\qquad
J_{22} = \int\Big[2\us p_2 + 2\mu\,\us_x (p_2)_x\Big]\dx .
\end{equation}

\subsection{Explicit Newton update}

With $\det\bm{J} = J_{11}J_{22} - J_{12}J_{21}$ and
\begin{equation}
\bm{J}^{-1} = \frac{1}{\det\bm{J}}
\begin{pmatrix} J_{22} & -J_{12} \\ -J_{21} & J_{11}\end{pmatrix},
\end{equation}
the update $\Delta\bm\lambda = -\bm{J}^{-1}\bm{F}$ is
\begin{equation}
\Delta\lambda_1 = -\frac{J_{22}F_1 - J_{12}F_2}{\det\bm{J}},
\qquad
\Delta\lambda_2 = -\frac{-J_{21}F_1 + J_{11}F_2}{\det\bm{J}} .
\end{equation}

\subsection{Compact discrete form}

Define the matrices
\begin{equation}
\bm{P} = \big[\,\bm{p}_1 \ \ \bm{p}_2\,\big] \in \R^{N\times2},
\qquad
\bm{R} = \big[\,\bm{r}_1 \ \ \bm{r}_2\,\big] \in \R^{N\times2},
\end{equation}
where $\bm{r}_j := \partial F_j/\partial\bm{\us}$ are the constraint gradients,
\begin{equation}
\bm{r}_1 = \bm{q},
\qquad
\bm{r}_2 = 2\,\bm{q}\odot\bm{\us} + 2\mu\,\bm{D}^\top\big(\bm{q}\odot\bm{D}\bm{\us}\big).
\label{eq:rvectors}
\end{equation}

\begin{lemma}[Factorisation of the Jacobian]
\label{lem:J}
$\bm{J} = \bm{R}^\top\bm{P}$.
\end{lemma}

\begin{proof}
By the chain rule,
$J_{ij} = \dfrac{\partial F_i}{\partial\lambda_j}
= \dfrac{\partial F_i}{\partial\bm{\us}}\cdot\dfrac{\partial\bm{\us}}{\partial\lambda_j}
= \bm{r}_i^\top\bm{p}_j = (\bm{R}^\top\bm{P})_{ij}$.
\end{proof}

\noindent
Lemma \ref{lem:J} is not merely cosmetic: it is what makes the backward pass in
Section \ref{sec:backward} collapse into a projector.

\subsection{Complete forward algorithm}

\begin{enumerate}[label=\textbf{\arabic*.}, leftmargin=2.2em]
\item Initialise $\lambda_1 = \lambda_2 = 0$ (or warm-start from the previous
      training step).
\item Set $\kappa = \lambda_2\mu/(1+\lambda_2)$ and assemble $\bm{A}$ as in
      \eqref{eq:matrixA}, respecting Remark \ref{rem:bcwarning}.
\item Build $\bm{h} = (\ut - \lambda_1/2)/(1+\lambda_2)$; overwrite
      $h_1 \leftarrow g_a$, $h_N \leftarrow g_b$.
\item Solve $\bm{A}\bm{\us} = \bm{h}$.
\item Compute $\bm{D}\bm{\us}$ and evaluate $F_1, F_2$ from \eqref{eq:discrete-F}.
\item If $|F_1| + |F_2| < \epsilon_{\text{tol}}$, go to step 10.
\item Solve $\bm{A}\bm{p}_1 = \bm{d}^{(1)}$ and
      $\bm{A}\bm{p}_2 = \bm{d}^{(2)}$, where the boundary rows are zero and,
      for $2\le i\le N-1$,
      \begin{equation}
      d^{(1)}_i = \frac{-1/2}{1+\lambda_2},
      \qquad
      d^{(2)}_i = \frac{\mu\,\delta_{xx}\us_i - \us_i}{1+\lambda_2} .
      \end{equation}
      The factor $1/(1+\lambda_2)$ is essential. The matrix $\bm{A}$
      discretises the \emph{normalised} operator
      $\mathcal{I}-\kappa\partial_x^2$, whereas \eqref{eq:p1} and
      \eqref{eq:p2} are stated for the unnormalised
      $\mathcal{H} = (1+\lambda_2)(\mathcal{I}-\kappa\partial_x^2)$. Dropping
      it inflates $\bm{P}$, and therefore $\bm{J}$, by a factor $1+\lambda_2$,
      corrupting both the Newton update and the backward projector. The
      stencil $\delta_{xx}$ must likewise be the same three-point operator
      that appears in $\bm{A}$; a wider stencil makes $\bm p_2$ inconsistent
      with the system and stalls the iteration.
\item Assemble $\bm{R}$ from \eqref{eq:rvectors} and $\bm{J} = \bm{R}^\top\bm{P}$.
\item Update $\bm\lambda \leftarrow \bm\lambda - \bm{J}^{-1}\bm{F}$; return to step 2.
\item Recompute $\bm{p}_1,\bm{p}_2,\bm{R},\bm{J}$ at the converged $\bm\lambda$
      and store them for the backward pass. Return $\bm{\us}$.
\end{enumerate}

\begin{remark}[Why step 10 matters]
Steps 7--8 are executed \emph{before} the convergence test of the following
iteration, so on exit the stored sensitivities correspond to the previous
iterate. They must be recomputed at the converged multipliers, otherwise the
gradient of Section \ref{sec:backward} is evaluated at the wrong point.
\end{remark}

\begin{remark}[Convergence and starting point]
Newton's method converges quadratically near the solution provided $\bm{J}$ is
non-singular, which holds whenever the constraint gradients $\bm{r}_1,\bm{r}_2$
are linearly independent, i.e.\ whenever the two constraints are not degenerate.

The starting point matters more than it might appear. From $\bm\lambda = 0$ the
iteration has to travel to $\lambda_2$ of order $30$ at late times, which needs
many damped steps and can stray onto the ill-posed branch $\lambda_2 < 0$.
Dropping the elliptic term gives the affine relation
$I_2[\us]\approx I_2[\ut]/(1+\lambda_2)^2$, which inverts to
\begin{equation}
\lambda_2^{(0)} = \sqrt{\frac{I_2[\ut]}{C_2}} - 1,
\qquad
\lambda_1^{(0)} = \frac{2}{|\Om|}\Big(\textstyle\int\ut\dx
- \big(1+\lambda_2^{(0)}\big)C_1\Big).
\end{equation}
This costs two quadratures and brings every time level to convergence within
about seven iterations. It is preferred here to warm-starting from the previous
training step: $\lambda_2$ varies strongly with $t$, from near zero at $t=0$ to
tens at late times, so a multiplier cached at one time level is a poor guess at
another when the levels are sampled at random. A backtracking line search
guards the remaining cases.
\end{remark}

%==========================================================================
\section{Backpropagation through the projection \textbf{[NEW]}}
\label{sec:backward}
%==========================================================================

Training requires the gradient of the loss with respect to the network
parameters. The projection sits between the network and the loss, so we need the
vector--Jacobian product
\begin{equation}
\bm{g}_{\ut} = \left(\frac{\partial\bm{\us}}{\partial\ut}\right)^{\!\top}\bm{g},
\qquad
\bm{g} := \frac{\partial\mathcal{L}}{\partial\bm{\us}} .
\end{equation}

\subsection{The two contributions}

The projected output depends on $\ut$ in two ways: directly through the
right-hand side of the linear solve, and indirectly because the multipliers
$\bm\lambda$ are themselves functions of $\ut$. By the chain rule,
\begin{equation}
\frac{\partial\bm{\us}}{\partial\ut}
= \underbrace{\frac{\partial\bm{\us}}{\partial\ut}\bigg|_{\bm\lambda}}_{\text{frozen multipliers}}
\;+\;
\underbrace{\bm{p}_1\left(\frac{\partial\lambda_1}{\partial\ut}\right)^{\!\top}
+ \bm{p}_2\left(\frac{\partial\lambda_2}{\partial\ut}\right)^{\!\top}}_{\text{multiplier response}} .
\label{eq:two-terms}
\end{equation}

\begin{remark}[This is where an implementation can go wrong]
Retaining only the first term of \eqref{eq:two-terms} gives a gradient whose
range is not tangent to the constraint manifold. The result is an approximate
descent direction rather than the exact implicit-function-theorem gradient.
The correction below costs one $2\times2$ solve and $O(N)$ work, so there is no
reason to omit it.
\end{remark}

\subsection{First term: frozen multipliers}

At fixed $\bm\lambda$, the system is $\bm{A}\bm{\us} = \bm{h}$ with
$\bm{h} = \bm{E}\,\ut/(1+\lambda_2) + \text{(boundary data)}$, where
$\bm{E} = \mathrm{diag}(0,1,\dots,1,0)$ zeroes the boundary rows (their values
$g_a, g_b$ do not depend on $\ut$). Differentiating,
\begin{equation}
\bm{A}\,\frac{\partial\bm{\us}}{\partial\ut}\bigg|_{\bm\lambda} = \frac{1}{1+\lambda_2}\bm{E}
\qquad\Longrightarrow\qquad
\frac{\partial\bm{\us}}{\partial\ut}\bigg|_{\bm\lambda}
= \frac{1}{1+\lambda_2}\,\bm{A}^{-1}\bm{E} =: \bm{M} .
\label{eq:M}
\end{equation}

\subsection{Second term: multiplier response by the implicit function theorem}

The multipliers satisfy $\bm{F}(\bm\lambda, \ut) = \bm{0}$ identically. Total
differentiation with respect to $\ut$ gives
\begin{equation}
\underbrace{\frac{\partial\bm{F}}{\partial\bm\lambda}}_{=\ \bm{J}}
\frac{\partial\bm\lambda}{\partial\ut}
+ \frac{\partial\bm{F}}{\partial\ut}\bigg|_{\bm\lambda} = \bm{0}
\qquad\Longrightarrow\qquad
\frac{\partial\bm\lambda}{\partial\ut}
= -\bm{J}^{-1}\frac{\partial\bm{F}}{\partial\ut}\bigg|_{\bm\lambda} .
\end{equation}
Now $F_i$ depends on $\ut$ only through $\bm{\us}$, so
\begin{equation}
\frac{\partial F_i}{\partial\ut}\bigg|_{\bm\lambda}
= \frac{\partial F_i}{\partial\bm{\us}}\,
  \frac{\partial\bm{\us}}{\partial\ut}\bigg|_{\bm\lambda}
= \bm{r}_i^\top\bm{M},
\qquad\text{i.e.}\qquad
\frac{\partial\bm{F}}{\partial\ut}\bigg|_{\bm\lambda} = \bm{R}^\top\bm{M} .
\end{equation}
Therefore
\begin{equation}
\frac{\partial\bm\lambda}{\partial\ut} = -\bm{J}^{-1}\bm{R}^\top\bm{M} .
\end{equation}

\subsection{Assembling the full Jacobian}

Substituting into \eqref{eq:two-terms} and writing the multiplier-response term
as $\bm{P}\,\partial\bm\lambda/\partial\ut$,
\begin{equation}
\frac{\partial\bm{\us}}{\partial\ut}
= \bm{M} - \bm{P}\bm{J}^{-1}\bm{R}^\top\bm{M}
= \Big(\bm{I} - \bm{P}\bm{J}^{-1}\bm{R}^\top\Big)\bm{M} .
\label{eq:full-jac}
\end{equation}

Using Lemma \ref{lem:J}, $\bm{J} = \bm{R}^\top\bm{P}$, so
\begin{equation}
\boxed{\
\frac{\partial\bm{\us}}{\partial\ut}
= \Big(\bm{I} - \bm{P}\big(\bm{R}^\top\bm{P}\big)^{-1}\bm{R}^\top\Big)\bm{M}
=: \big(\bm{I} - \bm\Pi\big)\bm{M} .
\ }
\end{equation}

\begin{proposition}[$\bm\Pi$ is an oblique projector]
$\bm\Pi := \bm{P}(\bm{R}^\top\bm{P})^{-1}\bm{R}^\top$ satisfies
$\bm\Pi^2 = \bm\Pi$. It projects onto $\mathrm{range}(\bm{P})$ along
$\mathrm{range}(\bm{R})^\perp$.
\end{proposition}

\begin{proof}
\begin{equation}
\bm\Pi^2 = \bm{P}(\bm{R}^\top\bm{P})^{-1}\underbrace{\bm{R}^\top\bm{P}}_{=\bm{J}}
(\bm{R}^\top\bm{P})^{-1}\bm{R}^\top
= \bm{P}(\bm{R}^\top\bm{P})^{-1}\bm{R}^\top = \bm\Pi . \qedhere
\end{equation}
\end{proof}

\begin{remark}[Interpretation]
The factor $\bm{I}-\bm\Pi$ removes exactly the components of the gradient that
would move the solution off the constraint manifold. The backward pass therefore
enforces conservation not only in the forward direction but also in the
optimisation direction: the network never receives a gradient signal telling it
to violate $C_1$ or $C_2$.
\end{remark}

\subsection{Practical backward algorithm}

Transposing \eqref{eq:full-jac},
\begin{equation}
\bm{g}_{\ut}
= \bm{M}^\top\Big(\bm{I} - \bm{R}\bm{J}^{-\top}\bm{P}^\top\Big)\bm{g}
= \bm{M}^\top\Big(\bm{g} - \bm{R}\,\bm{J}^{-\top}\bm{P}^\top\bm{g}\Big),
\end{equation}
which is evaluated in four steps:

\begin{enumerate}[label=\textbf{\arabic*.}, leftmargin=2.2em]
\item Form $\bm{s} = \bm{P}^\top\bm{g}
      = \big(\bm{p}_1^\top\bm{g},\ \bm{p}_2^\top\bm{g}\big)^\top \in \R^2$.
\item Solve the $2\times2$ system $\bm{J}^\top\bm{z} = \bm{s}$.
\item Form $\bm{g}_{\text{eff}} = \bm{g} - z_1\bm{r}_1 - z_2\bm{r}_2$.
\item Solve $\bm{A}^\top\bm{w} = \bm{g}_{\text{eff}}$ and set
      $\bm{g}_{\ut} = \bm{E}\bm{w}/(1+\lambda_2)$.
\end{enumerate}

\begin{remark}[Transpose, not inverse]
$\bm{A}$ in \eqref{eq:matrixA} is symmetric in its interior but \emph{not}
overall, because the boundary rows are $(1,0,\dots)$ and $(\dots,0,1)$ while the
adjacent columns retain $-c$. Step 4 must therefore solve with the genuine
transpose. Taking $(\bm{A}^\top)_{ij} = \bm{A}_{ji}$, the two zero entries move
diagonally:
\begin{equation}
(\bm{A}^\top)_{1,2} = -c,\quad (\bm{A}^\top)_{2,1} = 0,
\qquad
(\bm{A}^\top)_{N,N-1} = -c,\quad (\bm{A}^\top)_{N-1,N} = 0 .
\end{equation}
So $\bm{A}^\top$ is tridiagonal with diagonal $(1, 1+2c,\dots,1+2c, 1)$ and
$-c$ in every off-diagonal position \emph{except} $(\bm{A}^\top)_{2,1}$ and
$(\bm{A}^\top)_{N-1,N}$, which are zero. Filling those in would make the matrix
symmetric and feed the boundary residuals back into the interior incorrectly.
The cost is unchanged. In practice the safe implementation transposes the
assembled matrix rather than rebuilding it from this description.
\end{remark}

The backward pass requires one tridiagonal solve, one $2\times2$ solve, and
$O(N)$ vector arithmetic. It is no more expensive than the naive frozen-$\lambda$
version, and it is exact.

%==========================================================================
\section{Time derivative of the projected field \textbf{[NEW]}}
\label{sec:tangent}
%==========================================================================

Everything so far is stated at a fixed time level, which is correct for a
spatial projection but not yet enough for a PDE residual. The RLW residual
$r = u_t + u_x + \varepsilon uu_x - \mu u_{xxt}$ and the KdV residual
$r = u_t + 6uu_x + u_{xxx}$ both need the time derivative of the
\emph{projected} field, and $\us$ depends on $t$ through $\ut$ in a way that
couples the whole grid.

\subsection{The tangent is the same operator as the backward pass}

Write the projection as a map $\us(t) = \mathcal{P}\big[\ut(\cdot,t)\big]$.
Differentiating in $t$ by the chain rule, and using \eqref{eq:full-jac},
\begin{equation}
\boxed{\
\frac{\mathrm d\bm{\us}}{\mathrm dt}
= \frac{\partial\bm{\us}}{\partial\ut}\,\frac{\mathrm d\ut}{\mathrm dt}
= \big(\bm{I}-\bm\Pi\big)\bm{M}\,\frac{\mathrm d\ut}{\mathrm dt}. \ }
\label{eq:tangent}
\end{equation}
This is exact, not an approximation: it is the ordinary chain rule with the
Jacobian evaluated at $\ut(\cdot,t)$. The backward pass applies
$\big[(\bm I-\bm\Pi)\bm M\big]^\top$; the tangent applies the operator itself.
One derivation therefore serves both directions.

\subsection{Obtaining the raw tangent}

The quantity $\mathrm d\ut/\mathrm dt$ is ordinary automatic differentiation.
Both the network $N(x,t;\theta)$ and the hard-BC layer \eqref{eq:hardbc} act
pointwise in $t$, so $\partial\ut_i/\partial t_j = 0$ for $i \neq j$ and a
single reverse-mode pass with $\bm 1$ as the seed returns the entire vector.
The coupling is introduced only by the projection, and \eqref{eq:tangent}
accounts for it.

\subsection{Moving boundary data}

If $g_a$ and $g_b$ depend on $t$, the tangent inherits their derivatives.
Differentiating $\bm{A}\bm{\us} = \bm{h}$ in $t$, the interior rows of the
right-hand side carry $\mathrm d\ut/\mathrm dt$ divided by $1+\lambda_2$, while
the two boundary rows carry $g_a'(t)$ and $g_b'(t)$ directly, since
$h_1 = g_a$ and $h_N = g_b$. Hence
\begin{equation}
\bm{M}\bm v \;=\; \bm{A}^{-1}\bm{\rho},
\qquad
\rho_1 = g_a'(t),\quad
\rho_N = g_b'(t),\quad
\rho_i = \frac{v_i}{1+\lambda_2}\ \ (2\le i\le N-1),
\end{equation}
and $\us_t$ follows from \eqref{eq:tangent} as before. Setting the boundary
rows of $\bm\rho$ to zero, which is correct only for static data, would force
$\us_t(a,t) = 0$ and contradict $\us(a,t) = g_a(t)$.

\subsection{Higher derivatives and cost}

Spatial derivatives of the tangent follow from the fixed difference matrices,
since $\partial_t$ and $\partial_x$ commute:
\begin{equation}
\big(\us_{xx}\big)_t = \bm{D}_2\,\frac{\mathrm d\bm{\us}}{\mathrm dt},
\end{equation}
so the RLW residual needs no additional solves. Evaluating
\eqref{eq:tangent} costs one solve with $\bm{A}$ (not $\bm{A}^\top$), one
$2\times2$ solve, and $O(N)$ vector arithmetic: the same as the backward pass.

\begin{remark}[Which dependence is differentiated through]
In \eqref{eq:tangent} the state $(\bm{A},\bm{P},\bm{R},\bm{J},\lambda_2)$ is
evaluated at the current $\ut$ and then held fixed when the residual is
differentiated with respect to $\theta$. The tangent itself is exact; what is
dropped is the second-order dependence of the Jacobian on $\ut$. This is the
usual frozen-Jacobian treatment, and it affects the parameter gradient only at
second order.
\end{remark}

\begin{remark}[Why this matters for the method]
Prior projections avoid this issue by construction. PINN-Proj's correction is
$\us = \ut + \varkappa(t)$ and the affine projection is
$\us = \alpha(t)\ut + \beta(t)$; in both, the $t$-dependence sits in one or two
scalars, so ordinary automatic differentiation handles $\partial_t\us$ with no
extra machinery. The Helmholtz projection is a global operator on the grid, so
the tangent has to be propagated explicitly. That is the price of handling
derivative-dependent invariants, and \eqref{eq:tangent} shows it is a small
one.
\end{remark}

%==========================================================================
\section{Cost analysis \textbf{[NEW]}}
%==========================================================================

\begin{table}[h]
\centering
\begin{tabular}{llll}
\toprule
Stage & Operation & Count & Cost \\
\midrule
Forward, per Newton step
 & assemble $\bm{A}$              & 1 & $O(N)$ \\
 & tridiagonal solve              & 3 \ ($\bm\us,\bm p_1,\bm p_2$) & $O(N)$ \\
 & finite-difference derivatives  & 3 & $O(N)$ \\
 & quadratures ($F_1,F_2,\bm{J}$) & 5 & $O(N)$ \\
 & $2\times2$ solve               & 1 & $O(1)$ \\
\midrule
Backward
 & $\bm{P}^\top\bm{g}$            & 1 & $O(N)$ \\
 & $2\times2$ solve               & 1 & $O(1)$ \\
 & rank-2 correction              & 1 & $O(N)$ \\
 & tridiagonal solve with $\bm{A}^\top$ & 1 & $O(N)$ \\
\bottomrule
\end{tabular}
\caption{Cost of the Helmholtz projection. Typical Newton counts are 3--5 from a
cold start and 1--2 when warm-started.}
\end{table}

\noindent
For comparison, a single forward pass of an $L$-layer network of width $W$ over
$N$ points costs $O(NLW^2)$. With $L=8$, $W=50$ this is $2\times10^4$ operations
per point, against roughly $10N$ for the projection. The projection is therefore
negligible in arithmetic; the only implementation risk is kernel-launch overhead
from the sequential Thomas recursion.

%==========================================================================
\section{Extension to two dimensions \textbf{[NEW]}}
%==========================================================================

\subsection{Polynomial invariants (KP)}

For $C_1 = \iint u\dx\dy$ and $C_2 = \iint u^2\dx\dy$, the derivation of
Section \ref{sec:prior} Case C carries over verbatim with $\int\cdot\dx$
replaced by $\iint\cdot\dx\dy$ and $|\Om|$ by $|\Om_2| = (b_1-a_1)(b_2-a_2)$:
\begin{equation}
\us(x,y) = \frac{\sigma_0}{\tilde\sigma}\Big(\ut(x,y) - \tilde m\Big) + m_0 .
\end{equation}
Nothing in the derivation used the dimension, which is why the framework is
dimension-independent for polynomial invariants.

\subsection{Derivative-dependent invariants in 2D}

For an invariant $\iint\big(u^2 + \mu|\nabla u|^2\big)\dx\dy$, repeating the
variation of Theorem \ref{thm:el} with
\begin{equation}
\delta\!\iint \mu|\nabla \us|^2 \dx\dy
= \iint 2\mu\,\nabla \us\cdot\nabla v \dx\dy
= \oint 2\mu\,v\,\partial_n \us \,\mathrm{d}s - \iint 2\mu\,v\,\nabla^2\us\dx\dy
\end{equation}
(the boundary integral vanishing for $v|_{\partial\Om_2}=0$) gives
\begin{equation}
\boxed{\ (1+\lambda_2)\us - \lambda_2\mu\,\nabla^2\us = \ut - \frac{\lambda_1}{2}. \ }
\end{equation}
This is the two-dimensional Helmholtz equation. The Newton system remains
$2\times2$ regardless of dimension; only the elliptic solve changes, from a
tridiagonal solve to a sparse direct solve ($O(N^{3/2})$), an FFT-based solve
($O(N\log N)$), or multigrid ($O(N)$).

%==========================================================================
\section{Feasibility of hard $I_3$ \textbf{[NEW analysis]}}
%==========================================================================

\subsection{KdV energy}

For $C_3 = \int\big(u^3 - \tfrac12 u_x^2\big)\dx$, the variation of each piece is
\begin{equation}
\delta\!\int (\us)^3\dx = \int 3(\us)^2 v\dx,
\qquad
\delta\!\int\!-\tfrac12(\us_x)^2\dx = -\int \us_x v_x\dx = \int \us_{xx}v\dx ,
\end{equation}
so adding $\lambda_3 G_3$ to the Lagrangian contributes
$\lambda_3\big(3(\us)^2 + \us_{xx}\big)$. For KdV the momentum
$C_2 = \int u^2\dx$ carries \emph{no} derivative, so the $\lambda_2$ constraint
contributes only $2\lambda_2\us$ and the stationarity condition is
\begin{equation}
2(\us-\ut) + \lambda_1 + 2\lambda_2\us
+ \lambda_3\Big(3(\us)^2 + \us_{xx}\Big) = 0 .
\end{equation}
There is no $\lambda_2\mu\us_{xx}$ term here; that belongs to RLW's $I_2$, not
to KdV. The term $3\lambda_3(\us)^2$ is nonlinear in $\us$, so the
Euler--Lagrange equation is a nonlinear ODE. No single linear solve suffices,
and Newton on the full $N$-dimensional system would be required.

\subsection{RLW energy}

For $I_3 = \int(\varepsilon u^3+3u^2)\dx$ the contribution is
$\lambda_3\big(3\varepsilon(\us)^2 + 6\us\big)$, which contains no derivative.
Taken \emph{in isolation} the stationarity condition would be a pointwise
quadratic in $\us$, solvable node by node.

That is not the CE-PINN situation, because $I_2$ is enforced at the same time
and supplies the elliptic term. The condition for hard $I_1$, $I_2$ and $I_3$
together is
\begin{equation}
2(\us-\ut) + \lambda_1 + 2\lambda_2\us - 2\lambda_2\mu\,\us_{xx}
+ \lambda_3\Big(3\varepsilon(\us)^2 + 6\us\Big) = 0 ,
\end{equation}
a \emph{semilinear elliptic} boundary value problem. The term
$-2\lambda_2\mu\us_{xx}$ couples neighbouring nodes, so there is no node-by-node
solution unless $\lambda_2 = 0$ or $\mu = 0$.

\subsection{Decision}

$I_3$ is enforced as a soft penalty in relative form,
\begin{equation}
\mathcal{L}_{\text{soft}} = \frac{1}{N_t}\sum_{k=1}^{N_t}
\left(\frac{I_3\big[\us(\cdot,t_k)\big] - C_3}{C_3}\right)^{\!2} ,
\end{equation}
so that the weight $w_3$ does not have to absorb the scale of $C_3$, which
differs by orders of magnitude between the test problems. The penalty is
applied during the Adam phase and switched off for L-BFGS, so the final
descent minimises the PDE residual alone; this follows common practice for
conservation penalties and is a training choice, not a consequence of the
derivation. Hard enforcement of $I_3$ is stated as future work. The contribution of this
paper is the hard enforcement of $I_1$ and the derivative-dependent $I_2$.

%==========================================================================
\section{How the residual is evaluated \textbf{[NEW]}}
%==========================================================================

A standard PINN evaluates the residual at scattered collocation points by
automatic differentiation. That does not compose with a projection, because
the projection is a global operation on a spatial grid: there is no way to
project at an isolated point $(x_i,t_i)$ without the whole profile at $t_i$.

The residual is therefore evaluated on the same grid the projection uses. Each
training step draws $N_t$ time levels; at each level the network is evaluated
on the full spatial grid, Layer 1 and Layer 2 are applied, and the residual is
formed on the projected field $\us$. Spatial derivatives come from
finite-difference matrices, and the time derivative from
Section~\ref{sec:tangent}.

Two families of difference operators are used, deliberately kept apart. The
projection uses second-order operators throughout: the three-point Laplacian
inside $\bm{A}$, and the matching $\bm{D}$ above inside the invariant. Changing
either would break the correspondence between the discrete constraint and the
discrete Euler--Lagrange equation. The PDE residual uses fourth-order interior
stencils for $u_x$ and $u_{xx}$, and a fourth-order stencil for $u_{xxx}$ in
the KdV case, because residual accuracy is limited by these and not by the
projection. The outermost few nodes, where the one-sided stencils are less
accurate, are excluded from the residual norm; the hard boundary layer already
pins the solution there.

%==========================================================================
\section{The complete CE-PINN architecture \textbf{[NEW]}}
%==========================================================================

\subsection{Forward pass}

\begin{equation}
(x,t)
\ \xrightarrow{\ \text{MLP}\ }\ N(x,t;\theta)
\ \xrightarrow{\ \text{Layer 1, Eq.~\eqref{eq:hardbc}}\ }\ \ut
\ \xrightarrow{\ \text{Layer 2, Eq.~\eqref{eq:helmholtz}}\ }\ \us
\end{equation}

Layer 1 makes the initial and boundary conditions exact. Layer 2 makes $C_1$ and
$C_2$ exact. Both hold for every $\theta$, at every training step, and at
inference.

\subsection{Loss}

Because the boundary conditions and the first two invariants are structural, the
loss reduces to
\begin{equation}
\boxed{\
\mathcal{L}(\theta)
= \frac{1}{N_r}\sum_{i=1}^{N_r}\Big|r\big[\us\big](x_i,t_i)\Big|^2
\;+\; w_3\,\mathcal{L}_{\text{soft}} \ }
\end{equation}
with no $\mathcal{L}_{\text{IC}}$, no $\mathcal{L}_{\text{BC}}$, and no penalty
for $C_1$ or $C_2$. The residual is evaluated on the \emph{projected} field
$\us$, not on $\ut$: this is what makes the conservation constraint part of the
model rather than a post-processing step.

\subsection{Precomputation}

Before training, evaluate the initial condition on the projection grid and
compute the reference invariants by the same quadrature rule used during
training:
\begin{equation}
C_1 = \bm{q}^\top\bm{f},
\qquad
C_2 = \bm{q}^\top\Big[\bm{f}\odot\bm{f} + \mu(\bm{D}\bm{f})\odot(\bm{D}\bm{f})\Big],
\qquad
C_3 = \bm{q}^\top\Big[\varepsilon\,\bm{f}^{\odot3} + 3\,\bm{f}\odot\bm{f}\Big]
\qquad\text{(RLW).}
\end{equation}
The third invariant is equation-specific. For KdV,
\begin{equation}
C_3 = \bm{q}^\top\Big[\bm{f}^{\odot3}
      - \tfrac12(\bm{D}\bm{f})\odot(\bm{D}\bm{f})\Big]
\qquad\text{(KdV)},
\end{equation}
and for KP the corresponding Hamiltonian additionally carries the nonlocal
term $\tfrac32(\partial_x^{-1}u_y)^2$, which needs an auxiliary variable
$v$ with $v_x = u_y$ and is not enforced here.
Using the same quadrature for the reference values and for the constraints
ensures the discrete constraints are exactly satisfiable.

%==========================================================================
\section{Summary: what is prior work and what is new}
%==========================================================================

\begin{table}[h]
\centering
\small
\begin{tabular}{p{5.4cm} p{2.6cm} p{2.4cm} p{2.9cm}}
\toprule
Capability & PINN-Proj & SDIFP & \textbf{CE-PINN} \\
\midrule
$\int u\dx$ (linear)                      & Yes & Yes & Yes \\
$\int u^2\dx$ (polynomial)                & Yes & Yes & Yes \\
$\int(u^2+\mu u_x^2)\dx$ (derivative)     & \textbf{No} & \textbf{No} & \textbf{Yes} \\
$\int(u^3-\tfrac12u_x^2)\dx$ (derivative) & \textbf{No} & \textbf{No} & Analysed; soft \\
Hard Dirichlet boundary conditions        & No (Neumann only) & No (shifts boundary) & \textbf{Yes} \\
Projection preserves BCs                  & --- & No & \textbf{Yes, structurally} \\
Exact gradient through the projection     & Closed form & Closed form & \textbf{IFT + adjoint}$^\dagger$ \\
Two-dimensional demonstration             & Advection only & High-$d$ linear & \textbf{KP} \\
\bottomrule
\end{tabular}
\caption{Capability comparison. The rows in bold are the contribution of this
work. $^\dagger$The gradient of $\us$ with respect to $\ut$ is exact, by the
implicit function theorem. The tangent map $\ut_t \mapsto \us_t$ is exact at
the current state, but the state is frozen when the residual is differentiated
with respect to $\theta$, so the end-to-end parameter gradient omits the
second-order dependence of the projector on $\ut$. ``Projection preserves
BCs'' refers to the Helmholtz projection, where
the Dirichlet data enters as part of the elliptic boundary value problem. The
affine projection does not preserve non-homogeneous Dirichlet data, and
preserves homogeneous data only when $\beta = 0$; that applies equally to the
polynomial path of the present method when it is used on KdV and KP.}
\end{table}

\paragraph{The mathematical contribution, in one sentence.}
Prior projections solve an \emph{algebraic} Euler--Lagrange equation, which is
why their solution is an affine rescaling of the network output; the presence of
a derivative term in the invariant turns the Euler--Lagrange equation into an
\emph{elliptic boundary value problem}, and solving that problem, together with
the $2\times2$ Newton system for its multipliers and the adjoint gradient it
requires, is what this work adds.

\paragraph{Relationship to prior work.}
By Proposition \ref{prop:limit}, setting $\lambda_2 = 0$ recovers PINN-Proj's
linear projection exactly, and letting $\mu\to0$ recovers the joint affine
projection of PINN-Proj and SDIFP. Existing methods are the degenerate case of
the present method in which the elliptic term is absent.

\end{document}
